\documentclass[12pt]{article} % definisce tipo di documento
%Default usa testo carattere dimensione 10 -> con quadre si mette numero punti del carattere
\usepackage{graphicx} % Required for inserting images -> si possono usare pacchetti esterni
% Pacchetti richiamati con backslash \ all'inizio, e si aprono e chiudono parentesi graffe
\usepackage{hyperref} %permette di cliccare su riferimento come se fosse un collegamento -> utile per sommario
\usepackage{lineno}
\linenumbers %mette i numeri delle righe
\usepackage{comment}

\title{Il mio primo documento in LaTeX}
\author{Mirco Gruppi}
%\date{}
%\date{} senza data usa quella del sistema

\begin{document} %sempre seguito da un \end

\maketitle

\tableofcontents %inserisce automaticamente il sommario -> con hyperref manda a capitolo cliccando
\listoffigures
%mette lista delle immagini

\section{Introduzione} \label{Intro}
/label Mette sezioni del documento

label per richiamare nome paragrafo

Con spazio mette in automatico l'indentazione

\noindent 
/noindent funzione per non mettere indentazione automatica

\smallskip 
/smallskip aggiunge spazio dopo indentazione

\bigskip 
/bigskip aggiunge spazio più grande

\section{Metodi}
In questa tesi è stata utilizzata l'eugazione di Newton, raccolta nell'equazione \ref{eq:Newton}.
\begin{equation}
    F=G \times \frac{m_1 \times m_2}{r^2}
    % \frac fa funzioni mettendo denominatore e numeratore tra parentesi
    % _ underscore mette il pedice, ^ potenza
    % \times mette il per "x"
    % \sqrt mette sotto radice, tra parentesi quadre il numero dell'esponente \sqrt[3] 
    \label{eq:Newton} %etichetta
\end{equation}
\begin{equation}
    F=G \times \frac{\sqrt{m_1 \times \sqrt[3]{m_2}}}{r^2}
\end{equation}

Come visto nella sezione \ref{Intro} 

\begin{itemize}
    \item analisi multivariata
    \item analisi multitemporale
\end{itemize}

\begin{enumerate}
    \item cacio
    \item pepe
\end{enumerate}

\section{Risultati}
Assurdo come sia come in figura \ref{fig:MARINOCAVALLO}

\begin{figure} %immagine spostata per evitare buchi vuoti
    \centering %centra immagini
    \includegraphics[width=0.5\linewidth]{Nuovo progetto (1).png} %width moltiplica larghezza
    \caption{WE LOVE MARINO CAVALLO}
    \label{fig:MARINOCAVALLO}
\end{figure}

\begin{figure}
    \centering
    \includegraphics[width=0.25\linewidth]{ezgif.com-webp-to-png-converter.png}
    \caption{Zazzingara}
    \label{fig:zazzingara}
\end{figure}

\section{Discussione}

%begin{comment}
    Se devo togliere una parte ma voglio conservarla, basta mettere sotto comment e ti toglie tutta una sezione del testo dal foglio prodotto, ma posso richiamarlo quando voglio
%end{comment} %si può riportare commentando end comment


\end{document}
